\chapter{Verifica e Validazione}
\label{chap:verifica_validazione}

\textit{Questo capitolo raccoglie le attività di verifica e validazione condotte sistema di
telemetria EggLog, con l'obiettivo di dimostrarne sia la correttezza funzionale
in condizioni d'uso reali, sia l'adeguatezza prestazionale rispetto ai vincoli
del sistema embedded su cui è eseguito.} 

\section{Introduzione}

Il capitolo è organizzato in due sezioni: la prima presenta i test svolti,
distinguendo tra i test su campo, condotti direttamente su motocicletta, e i
test su banco, pensati per verificare in modo controllato singoli
comportamenti del firmware; la seconda sezione presenta i \textit{benchmark}
prestazionali eseguiti sul firmware, relativi al tempo di scrittura su
microSD e al carico computazionale imposto al microcontrollore.

\section{Test}
\label{sec:test}

\subsection{Metodologia dei test su campo}
\label{subsec:metodologia_test_campo}

I test su campo sono stati condotti avviando il firmware in modalità
\textit{standalone}, così da riprodurre le condizioni d'uso reali del
dispositivo: il boot avviene automaticamente non appena il sistema
entra in contatto con una sorgente di alimentazione, senza alcun intervento
manuale aggiuntivo.

\begin{enumerate}
    \item \textbf{Allestimento del prototipo.} Per ciascun test è stato
    utilizzato il prototipo saldato (non la versione su breadboard, che
    ovviamente non avrebbe retto alle sollecitazioni di un utilizzo su
    strada). Il prototipo è stato fissato con dello scotch al serbatoio di
    una Yamaha MT-07. Sotto il codino della moto è stato collegato alla
    porta diagnostica l'EggLog CAN Sniffer, alimentato tramite un
    powerbank; analogamente, anche EggLog è stato collegato a un
    powerbank, riposto nella tasca del pilota.

    \item \textbf{Preparazione della sessione.} Prima di ogni sessione di
    test è stata eseguita una pulizia completa della microSD, in modo da
    riportare il dispositivo in una condizione equivalente a quella di
    \textit{appena acquistato}. È stato inoltre fissato al manubrio un
    dispositivo mobile con installata l'applicazione EggLog App, collegata
    a EggLog, ed è stata eseguita la calibrazione dell'IMU per configurare
    correttamente tutti gli assi. A questo punto è stata avviata una
    sessione di guida normale su strada, della durata complessiva di
    15/20 minuti.

    \item \textbf{Raccolta e conversione dei dati.} Al termine di ciascuna
    sessione, il sistema è stato smontato, la microSD estratta e i file
    binari recuperati. Questi ultimi sono stati convertiti in formato CSV
    tramite uno script Python sviluppato ad hoc, capace di leggere tutti i
    file \texttt{.bin} generati e produrre i corrispondenti file
    \texttt{.csv} per la successiva analisi.
\end{enumerate}

\vspace{0.5cm}

L'analisi dei dati raccolti è stata condotta incrociando il percorso
realmente effettuato (e i relativi timestamp) con l'insieme dei dati
registrati, al fine di individuare eventuali bug, buchi nei dati, e per
verificare che i valori registrati fossero coerenti con quanto effettivamente
accaduto in un dato istante. \\

Alcuni esempi dei criteri di verosimiglianza adottati sono:

\begin{itemize}
    \item in ingresso a una rotonda, l'accelerazione frontale deve essere
    negativa, la velocità deve calare e l'inclinazione frontale deve
    essere negativa (il posteriore si solleva leggermente);

    \item durante una curva a sinistra percorsa a velocità costante,
    l'accelerazione frontale deve essere prossima a zero, mentre
    l'inclinazione laterale deve avere segno negativo e rientrare in un range
    verosimile;

    \item in fase di accelerazione, l'accelerazione frontale deve essere
    positiva, così come l'inclinazione frontale, segno che la moto tende
    a impennarsi leggermente per effetto dell'allungamento delle forcelle.
\end{itemize}

\label{subsec:test_campo}
Sono stati condotti tre test su campo, secondo la metodologia appena
descritta. I primi due hanno portato all'individuazione di alcune
criticità, poi risolte; il terzo si è invece concluso senza criticità, 
confermando il raggiungimento degli obiettivi prefissati.


\subsection{Test su campo 1 - Calibrazione e posizionamento del sensore}
\label{subsubsec:test_campo_1}

Il primo test ha evidenziato due criticità distinte, accennate anche in
\hyperref[subsec:prob_risc_calib_imu]{Problematiche riscontrate durante la
calibrazione dell'IMU}.

Il primo step della calibrazione era caratterizzato da una soglia di varianza
troppo restrittiva, che non teneva conto delle vibrazioni trasmesse dal
motore acceso. Di conseguenza, con la moto accesa, il tremolio impediva il
completamento con successo del primo passaggio di calibrazione. Il problema
è stato temporaneamente risolto eseguendo la calibrazione a moto spenta.

Il giro di prova è stato completato con successo e la moto 3D
dell'applicazione ha replicato fedelmente i movimenti reali del veicolo. È
emerso tuttavia un secondo aspetto: essendo il sensore posizionato proprio
sopra il serbatoio, l'accelerazione risultava accentuata, venendo talvolta
interpretata dal firmware come un'impennata più marcata di quanto realmente
avvenuto. La spiegazione, coerente con la posizione di montaggio, è che il
sensore, trovandosi a una certa distanza dal centro di rotazione del
veicolo, risenta di un effetto di braccio di leva: un lieve sollevamento del
frontale, nel punto in cui è montato il prototipo, si traduce in
un'inclinazione percepita superiore a quella reale.

In questa fase, il confronto con i dati provenienti dallo sniffer del bus CAN
si è rivelato particolarmente utile: trattandosi di dati più affidabili
rispetto a quelli restituiti dai sensori del prototipo, hanno permesso di
effettuare l'incrocio con maggiore semplicità e sicurezza.


\subsection{Test su campo 2 - Arrotondamento dei dati di accelerazione}
\label{subsubsec:test_campo_2}

Nel secondo test, la fase di calibrazione (già corretta a seguito del test
precedente) è andata a buon fine senza criticità. Un'analisi più
approfondita dei dati raccolti ha però permesso di individuare un problema
non emerso a una prima ispezione: i valori di accelerazione risultavano
sistematicamente arrotondati. La causa è stata ricondotta a un errore di
conversione nel firmware, introdotto nel momento in cui i dati venivano
inseriti nel superfile (\hyperref[sec:impl-superfile]{Sistema di persistenza
a superfile}).

Il problema è stato corretto e i dati rimanenti sono stati analizzati
comunque, confermando la correttezza del comportamento del sistema.

\subsection{Test su campo 3 - Validazione finale}
\label{subsubsec:test_campo_3}

Il terzo test si è svolto senza che emergessero criticità. L'analisi
incrociata dei dati raccolti con il percorso effettuato non ha evidenziato
bug, buchi nei dati o incoerenze rispetto alla dinamica di guida attesa. Il
tutor aziendale ha confermato, sulla base di questi risultati, il
raggiungimento dell'obiettivo finale.


\subsection{Test su banco 1 - Sospensione e ripresa del logging con fix GPS}
\label{subsubsec:test_banco_1}
\label{subsec:test_banco}

Oltre ai test su campo, sono stati condotti due test su banco, con
l'obiettivo di verificare in modo controllato la gestione dell'inattività
dei sensori (\hyperref[sec:inattivita]{Gestione dell'inattività dei
sensori}) e il comportamento del firmware in assenza di movimento, sia con
che senza fix GPS.

Per verificare la gestione dell'inattività dei sensori, la procedura seguita
è stata la seguente:

\begin{enumerate}
    \item il firmware viene temporaneamente modificato per impostare a 5
    minuti il timer minimo di inattività configurabile;
    \item il dispositivo viene portato all'aperto, in attesa
    dell'acquisizione del fix GPS;
    \item avviata una sessione di logging, il dispositivo viene lasciato
    perfettamente immobile;
    \item trascorsi i 5 minuti, il LED della ESP32 inizia a pulsare in 
    verde, a indicare l'ingresso nello stato di sospensione: il test ha 
    quindi avuto esito positivo;
    \item il dispositivo viene ripreso in mano e spostato: il LED torna
    verde fisso, segnalando il rientro nello stato di logging e la ripresa
    della scrittura dei dati.
\end{enumerate}


\subsection{Test su banco 2 - \textit{Fallback} sull'IMU senza fix GPS}
\label{subsubsec:test_banco_2}

Lo stesso procedimento è stato ripetuto senza fix GPS, direttamente sulla
scrivania in ufficio, per verificare il comportamento di \textit{fallback}
basato sull'IMU. Il risultato è stato identico a quello del test precedente:
sospensione dopo il timeout configurato e ripresa del logging non appena il
dispositivo veniva mosso.

\subsection{Sintesi dei test}
\label{subsec:sintesi_test}

I test sono identificati mediante un codice univoco composto da due elementi:
\texttt{X-N}, dove \texttt{X} identifica la modalità di esecuzione del test e 
\texttt{N} rappresenta il numero progressivo del test. In particolare, \texttt{C} 
indica un test condotto su campo, mentre \texttt{B} indica un test condotto su banco.

\begin{table}[H]
    \centering
    \begin{tabular}{p{2.0cm} p{2.8cm} p{8cm}}
        \hline
        \textbf{Codice} &
        \textbf{Esito} &
        \textbf{Scoperta / Nota} \\
        \hline

        \texttt{C-1} &
        Parzialmente positivo &
        Soglia di varianza della calibrazione troppo restrittiva; effetto di accentuazione
        dell'accelerazione per posizionamento sopra il serbatoio \\

        \texttt{C-2} &
        Parzialmente positivo &
        Arrotondamento dei dati di accelerazione per errore di conversione
        nel superfile \\

        \texttt{C-3} &
        Positivo &
        Nessuna criticità; conferma del raggiungimento degli obiettivi \\

        \texttt{B-1} &
        Positivo &
        Sospensione e ripresa del logging corrette, con fix GPS \\

        \texttt{B-2} &
        Positivo &
        Sospensione e ripresa del logging corrette, in fallback su IMU senza
        fix GPS \\

        \hline
    \end{tabular}
    \caption{Sintesi dei test condotti sul firmware}
    \label{tab:sintesi_test}
\end{table}

\newpage

\section{Benchmark prestazionali}
\label{sec:benchmark_prestazionali}

\subsection{Metodologia}
\label{subsec:metodologia_benchmark}

I benchmark relativi al tempo di scrittura sono stati realizzati introducendo un timer 
attorno a ciascuna operazione che
prevedeva l'accesso al modulo di scrittura su microSD. Per ogni operazione veniva
avviato immediatamente prima dell'inizio dell'accesso al modulo e
arrestato al completamento della relativa procedura di scrittura.

Il tempo così misurato veniva quindi passato a un oggetto statico dedicato alla
raccolta delle statistiche. Attraverso metodi specifici, tale oggetto manteneva
e aggiornava progressivamente il tempo minimo, medio e massimo di scrittura osservati 
durante l'esecuzione del benchmark. \\

Le operazioni sottoposte a misurazione sono state progettate appositamente per
avere lo stesso peso computazionale e le stesse caratteristiche in entrambe le
versioni del firmware. Ciò ha permesso di applicare la medesima metodologia sia
all'implementazione precedente, basata sulla gestione di un file per ciascun
\textit{chunk}, sia alla nuova implementazione basata sul superfile, rendendo i risultati
direttamente confrontabili.


\subsection{Tempo di scrittura su microSD}
\label{subsec:tempo_scrittura_sd}

\paragraph{Il problema della cluster size (superato)}
\label{subsubsec:problema_cluster_size}
Nell'approccio precedente, basato su un file per \textit{chunk} 
(\hyperref[sec:situazione-partenza]{Situazione di partenza}), il
dimensionamento della cluster size rappresentava un problema di non semplice
soluzione: bisognava gestire il rischio di frammentazione. Ogni nuovo file creato
richiedeva infatti l'allocazione di nuovi cluster, con un relativo \textit{overhang}
ogni volta che il \textit{chunk} non riempiva completamente l'ultimo cluster
allocato.

Con l'introduzione del superfile (\hyperref[sec:superfile]{Sistema di
persistenza a superfile}), un unico file preallocato per sensore, mai chiuso
né ricreato per l'intera durata del logging, ha eliminato alla radice questo
problema: l'area dati viene allocata una sola volta per l'intera retention,
e tutte le scritture successive avvengono all'interno di cluster già
assegnati. Non è quindi più necessario gestire l'overhang \textit{chunk} per \textit{chunk},
né si introduce frammentazione a runtime.


\paragraph{Confronto misurato: un file per chunk vs superfile}
\label{subsubsec:confronto_scrittura}
Per quantificare il beneficio introdotto dal superfile, è stato misurato il
tempo (in $\mu\mathrm{s}$) dell'intero blocco di scrittura su un campione di 8 
ore di logging:

\begin{table}[H]
    \centering
    \begin{tabular}{lrrr}
        \hline
        &
        \textbf{Un file per chunk} &
        \textbf{superfile} &
        \textbf{Differenza} \\
        \hline

        avg &
        $150.045\,\mu\mathrm{s}$ &
        $83.330\,\mu\mathrm{s}$ &
        \textbf{$-$44\%} \\

        min &
        $70.547\,\mu\mathrm{s}$ &
        $53.447\,\mu\mathrm{s}$ &
        $-$24\% \\

        max &
        $665.726\,\mu\mathrm{s}$ &
        $131.747\,\mu\mathrm{s}$ &
        \textbf{$-$80\%} \\

        \hline
    \end{tabular}
    \caption{Confronto dei tempi di scrittura su microSD}
    \label{tab:benchmark_scrittura_sd}
\end{table}

I risultati mostrano un netto miglioramento su tutte le metriche
considerate, in particolare sul valore massimo, dove il superfile riduce
dell'80\% il caso peggiore osservato. È inoltre da notare che, aumentando la
dimensione del campione, il tempo medio della versione di partenza tende a
peggiorare ulteriormente, per effetto della crescente frammentazione e del
tempo di ricerca dello spazio libero, mentre quello del superfile rimane
stabile: un comportamento coerente con l'assenza, in quest'ultimo, di
allocazione di nuovi cluster a runtime.


\subsection{Occupazione di CPU}
\label{subsubsec:cpu-occupazione}
\label{subsec:cpu_ram}

Oltre al tempo di scrittura su microSD, sono stati misurati anche
l'occupazione di CPU nei diversi stati del firmware e l'occupazione di RAM
effettivamente utilizzata dai buffer dei sensori, per entrambe le versioni.

Poiché il logging veniva avviato automaticamente non appena il dispositivo
veniva collegato a una fonte di alimentazione (\hyperref[sec:situazione-partenza]{Situazione
di partenza}), non esisteva uno stato di attesa distinto dallo stato di
logging: l'occupazione di CPU restava quindi costantemente attorno al
25\%.

Con l'introduzione del protocollo di avvio e arresto della sessione
(\hyperref[sec:avvio-arresto]{Protocollo di avvio e arresto della
sessione}), il dispositivo, appena collegato all'alimentazione e senza
alcuna sessione avviata, presenta un'occupazione di CPU di circa il 3\%.
All'avvio effettivo di una sessione, il consumo sale a circa il 23\%, un
valore inferiore anche a quello della versione di partenza. L'introduzione
della task dedicata alla segnalazione tramite LED
(\hyperref[sec:led]{Segnalazione di stato tramite LED}), assente nella
versione originale, aggiunge un piccolo \textit{overhead} rispetto alle sole
task di producer e consumer già presenti; tale \textit{overhead} risulta però
più che compensato dalla configurazione attuale dei sensori, in cui il
magnetometro è disabilitato: una task producer in meno comporta un risparmio
di cicli di CPU superiore al costo della nuova task LED. Il firmware attuale
tende inoltre a risparmiare più cicli di CPU possibile rispetto alla
versione di partenza grazie all'\textit{inactivity} dei sensori implementata.

La Tabella \ref{tab:cpu-occupazione} riassume i valori osservati nei diversi
stati.

\begin{table}[H]
    \centering
    \begin{tabular}{p{7.5cm} c c}
        \hline
        \textbf{Stato del dispositivo} & \textbf{Prima} & \textbf{Dopo} \\
        \hline
        Prima di una sessione & 25\% & 3\% \\
        Durante una sessione & 25\% & 23\% \\
        \hline
    \end{tabular}
    \caption{Occupazione di CPU nei diversi stati del firmware}
    \label{tab:cpu-occupazione}
\end{table}


\subsection{Occupazione della RAM}
\label{subsubsec:ram-occupazione}

\paragraph{Calcolo della RAM occupata}
Entrambe le versioni assegnano a ogni producer due buffer alternati
(\textit{ping-pong}, si veda \hyperref[sec:situazione-partenza]{Situazione
di partenza}), quindi la RAM effettivamente occupata da ciascun
sensore è pari al doppio della dimensione del singolo buffer.

\paragraph{Confronto per sensore}
La Tabella \ref{tab:ram-buffer-sensori} riporta il confronto tra le due
versioni, entrambe i buffer inclusi. Gli aumenti più
significativi riguardano i sensori GPS e IMU: la relativa giustificazione è
approfondita nel paragrafo seguente. Il buffer del barometro resta invece
sostanzialmente invariato. Il magnetometro, il cui buffer nella versione
attuale sarebbe stato dimensionato per circa 31.8 KB, risulta invece
disabilitato: l'occupazione di RAM ad esso associata è quindi nulla (si veda
il paragrafo successivo).

\begin{table}[H]
    \centering
    \begin{tabular}{l r r r}
        \hline
        \textbf{Sensore} & \textbf{Versione di partenza} & \textbf{Versione attuale} & \textbf{Differenza} \\
        \hline
        Barometro & 5.0 KB & 4.7 KB & $-$0.3 KB \\
        GPS & 4.0 KB & 39.4 KB & +35.4 KB \\
        IMU & 123.0 KB & 130.3 KB & +7.3 KB \\
        Magnetometro & 32.0 KB & Disabilitato & $-$32.0 KB \\
        \hline
        \textbf{Totale} & \textbf{164.0 KB} & \textbf{174.4 KB} & \textbf{+10.4 KB} \\
        \hline
    \end{tabular}
    \caption{RAM occupata dai buffer dei producer (ping-pong incluso)}
    \label{tab:ram-buffer-sensori}
\end{table}

\paragraph{Variazione percentuale tra le due versioni}
Espresse in termini percentuali rispetto alla versione di partenza, le
variazioni per ciascun sensore sono riportate in
Tabella \ref{tab:ram-buffer-variazione-percentuale}.

\begin{table}[H]
    \centering
    \begin{tabular}{l r}
        \hline
        \textbf{Sensore} & \textbf{Variazione} \\
        \hline
        Barometro & $-$6.0\% \\
        GPS & +885.0\% \\
        IMU & +5.9\% \\
        Magnetometro & Disabilitato \\
        \hline
        \textbf{Totale} & \textbf{+6.3\%} \\
        \hline
    \end{tabular}
    \caption{Variazione percentuale della RAM occupata dai buffer dei producer}
    \label{tab:ram-buffer-variazione-percentuale}
\end{table}

\paragraph{Il ruolo del formato di serializzazione}
Il confronto in termini di \textit{kilobyte} occupati, per quanto corretto, non
racconta però l'intera storia: nella versione di partenza ogni campione
veniva serializzato in formato testuale CSV direttamente all'interno del
buffer del producer, mentre nella versione attuale ogni campione occupa
uno slot a dimensione fissa pari a \texttt{RECORD\_SIZE} in un formato
binario compatto. Stimando la lunghezza di una riga CSV a partire dai
formati di conversione impiegati (\texttt{snprintf}) e dai relativi range
di valore dei singoli campi, si ottiene una dimensione per campione
sensibilmente superiore a quella del corrispondente record binario, come
riportato in Tabella \ref{tab:csv-vs-binary-record}.

\begin{table}[H]
    \centering
    \begin{tabular}{l c c}
        \hline
        \textbf{Sensore} & \textbf{CSV (B/record)} & \textbf{Binario (B/record)} \\
        \hline
        Barometro & 32 & 20 \\
        GPS & 75 & 42 \\
        IMU & 56 & 20 \\
        Magnetometro & 40 & 21 \\
        \hline
    \end{tabular}
    \caption{Confronto tra la dimensione stimata di un record in formato CSV e la dimensione del record binario attuale}
    \label{tab:csv-vs-binary-record}
\end{table}

Applicando queste dimensioni alla capacità del singolo buffer riportata in
Tabella \ref{tab:ram-buffer-sensori}, emerge che il numero di campioni
effettivamente contenibili in un buffer è cresciuto molto più di quanto
suggerito dal solo aumento in kilobyte, ed è aumentato persino per il
barometro, il cui buffer si è invece \textit{ridotto} in termini assoluti
(Tabella \ref{tab:ram-buffer-variazione-percentuale}).

\begin{table}[H]
    \centering
    \begin{tabular}{l c c c}
        \hline
        \textbf{Sensore} & \textbf{Record/buffer (prima)} & \textbf{Record/buffer (ora)} & \textbf{Fattore} \\
        \hline
        Barometro & $\approx 80$ & 120 & $\times 1{,}5$ \\
        GPS & $\approx 27$ & 480 & $\times 17{,}8$ \\
        IMU & $\approx 1124$ & 3336 & $\times 3{,}0$ \\
        \hline
    \end{tabular}
    \caption{Numero stimato di campioni contenuti in un singolo buffer, prima e dopo il passaggio al formato binario}
    \label{tab:record-per-buffer}
\end{table}

In altre parole, l'incremento di RAM riportato in
Tabella \ref{tab:ram-buffer-variazione-percentuale} rappresenta solo una
parte del guadagno complessivo: una quota rilevante dell'aumento di
capacità di ritenzione dei buffer deriva dal passaggio da una
rappresentazione testuale, verbosa e a lunghezza variabile, a una
rappresentazione binaria compatta a lunghezza fissa, indipendente
dall'eventuale ridimensionamento dei buffer stessi.

\paragraph{Un compromesso tra RAM e frequenza di scrittura}
L'aumento dei buffer di GPS e IMU deriva dalla scelta di privilegiare un minor
numero di operazioni di I/O sulla microSD, a fronte di un maggiore utilizzo
della RAM. Le operazioni di I/O sono infatti sensibilmente più lente rispetto
all'accesso alla memoria principale; accumulare più campioni prima di ogni
scrittura consente quindi di ridurne la frequenza e il relativo costo
prestazionale. Tale dimensionamento è indipendente dalla finestra di retention
del superfile
(\hyperref[subsec:capacita_superfile]{Capacità del file}) e dipende dalla
frequenza di trasferimento dei campioni
(\texttt{FLUSH\_INTERVAL\_S}, Tabella \ref{tab:constants-per-sensore}).

La scelta è stata applicata a GPS e IMU, considerati più critici rispetto al
barometro, per i quali si è preferito sostenere una maggiore occupazione della
RAM pur di ridurre gli accessi alla microSD. Per evitare la perdita dei campioni
ancora presenti nei buffer al termine della sessione, la sequenza di arresto
(\hyperref[sec:avvio-arresto]{Protocollo di avvio e arresto della sessione})
esegue uno svuotamento forzato di tutti i buffer prima della chiusura dei
superfile.

\paragraph{Disabilitazione del magnetometro}
Il magnetometro, i cui dati non trovano al momento alcun utilizzo concreto
nel sistema, è stato invece disabilitato in questa versione del firmware: il
suo buffer, presente nella versione di partenza con una dimensione di circa
32.0 KB, non viene più allocato, azzerandone completamente l'occupazione di
RAM.

\paragraph{Il momento dell'allocazione conta più della dimensione}
Il dato più rilevante non riguarda quindi la dimensione dei singoli buffer,
quanto il momento in cui vengono allocati. Nella versione di partenza, il
logging iniziava automaticamente all'accensione e tutti e quattro i buffer
restavano allocati permanentemente, per un totale di circa 164 KB. Nella
versione attuale, invece, i buffer dei sensori attivi (barometro, GPS e IMU)
vengono allocati solo all'avvio effettivo di una sessione, grazie al protocollo
di avvio (\hyperref[sec:avvio-arresto]{Protocollo di avvio e arresto della
sessione}) e ai producer inattivi fino alla ricezione di un comando esplicito.
Nello stato di attesa risultano quindi liberi circa 174 KB, contribuendo al
netto calo dell'occupazione di CPU, dal 25\% al 3\%
(Tabella \ref{tab:cpu-occupazione}).

Durante il logging, l'occupazione dei buffer aumenta invece a circa 174 KB,
contro i 164 KB della versione di partenza. L'incremento, pari a circa il 6\%,
è un compromesso deliberato: si è preferito utilizzare una maggiore quantità
di RAM per accumulare più campioni e ridurre la frequenza delle operazioni di
I/O sulla microSD, più lente rispetto all'accesso alla memoria principale.

\newpage